\chapitrepage{Présentation du cadre du projet}

\begin{objectifbox}
Ce chapitre situe le projet dans son environnement professionnel. Nous présentons
le domaine d'activité, le processus de décision avant informatisation, les limites
observées, la solution retenue, son originalité et l'organisation des travaux.
\end{objectifbox}

\section{Introduction}

La valeur d'un système de prévision ne dépend pas uniquement de la qualité d'un
algorithme. Elle dépend également de la fraîcheur des données, de la manière dont
les résultats sont expliqués, du lien avec les contraintes de production et de la
capacité des utilisateurs à corriger une situation exceptionnelle. Nous avons
donc commencé par comprendre le fonctionnement du point de vente et les décisions
qui devaient être améliorées avant de choisir les technologies.

\section{Présentation de l'organisme d'accueil}

\subsection{Contexte général}

PAUL est une enseigne de boulangerie et de restauration dont l'identité de marque
rappelle une maison fondée en 1889. Le contexte étudié concerne un point de vente
au Maroc. Son activité combine des produits de boulangerie, de viennoiserie, de
pâtisserie, de cuisine et de boissons. Cette diversité crée des profils de demande
très différents : certains produits sont vendus chaque jour en volume élevé,
d'autres sont intermittents, saisonniers ou associés à un moment de consommation
précis.

Le projet se concentre sur l'interface entre trois fonctions opérationnelles :

\begin{itemize}
  \item la \textbf{vente}, qui fournit les tickets et les quantités réellement
  écoulées ;
  \item la \textbf{production}, qui doit préparer les bonnes quantités au bon
  moment ;
  \item l'\textbf{approvisionnement}, qui doit convertir le plan de production en
  besoins de farine, beurre, lait, fruits, emballages et autres composants.
\end{itemize}

Les équipes manipulent également des informations transversales : fêtes
marocaines, matchs de l'équipe nationale, événements ponctuels, commandes de
clients professionnels, recettes validées par le chef et prix d'achat. La solution
devait rassembler ces éléments sans déplacer les données sensibles vers un service
externe.

\subsection{Périmètre métier}

Le corpus local audité contient 373\,875 lignes de ventes journalières, comprises
entre le 7 juillet 2021 et le 28 juin 2026. Il couvre 1\,131 libellés produits
répartis dans dix familles et représente 4\,364\,539 unités vendues. Ces ordres de
grandeur imposent une automatisation : un traitement manuel produit par produit
serait trop lent et difficilement reproductible.

\begin{table}[H]
\centering
\caption{Périmètre de données observé dans la version auditée}
\label{tab:perimetre-donnees}
\begin{tabularx}{\textwidth}{>{\bfseries}p{4.4cm}Y}
\toprule
Indicateur & Valeur observée \\
\midrule
Période des ventes & 07/07/2021 au 28/06/2026 \\
Lignes de ventes journalières & 373\,875 \\
Jours distincts & 1\,817 \\
Produits distincts & 1\,131 \\
Familles de produits & 10 \\
Quantité totale enregistrée & 4\,364\,539 unités \\
Chiffre d'affaires TTC cumulé & 57,36 millions MAD \\
\bottomrule
\end{tabularx}
\end{table}

\section{Étude de l'existant}

\subsection{Description du processus existant}

Les ventes historiques proviennent d'abord d'états générés par le logiciel de
caisse PI Électronique/Elyx. Une partie de l'historique est disponible sous forme
d'exports DOS à largeur fixe et de classeurs mensuels. Les ventes les plus récentes
peuvent être extraites depuis SQL Server. Par ailleurs, les recettes techniques et
les prix matières sont conservés dans des fichiers distincts. Les décisions de
production reposent donc sur plusieurs supports, dont les structures, les noms de
produits et les granularités ne sont pas uniformes.

Avant l'intégration proposée, le processus pouvait être résumé comme suit :

\begin{enumerate}
  \item consulter les ventes passées ou les états de caisse ;
  \item estimer les quantités futures à partir de l'expérience récente ;
  \item ajuster mentalement cette estimation selon le calendrier ;
  \item reporter les quantités dans des supports opérationnels ;
  \item rapprocher manuellement les produits de leurs matières premières.
\end{enumerate}

Cette organisation contient une connaissance métier utile, mais elle rend la
décision dépendante de la disponibilité d'une personne, ne mesure pas
systématiquement l'erreur et complique l'explication d'une recommandation.

\subsection{Critique de l'existant}

L'analyse a mis en évidence les insuffisances suivantes :

\begin{itemize}
  \item \textbf{dispersion des données} : ventes mensuelles, ventes journalières,
  recettes, événements et prix ne sont pas réunis dans une chaîne unique ;
  \item \textbf{hétérogénéité des sources} : encodage CP850, sous-totaux présents
  dans certains exports et variations de libellés peuvent introduire des doubles
  comptes ou des ruptures de séries ;
  \item \textbf{prévision uniforme} : une même méthode ne convient ni à un produit
  régulier, ni à une vente intermittente, ni à un produit récemment introduit ;
  \item \textbf{absence de validation temporelle systématique} : une précision
  calculée sur les données d'entraînement surestime la performance réelle ;
  \item \textbf{prise en compte informelle des événements} : les fêtes, matchs et
  commandes importantes ne sont pas intégrés de manière reproductible ;
  \item \textbf{rupture entre vente et achat} : une prévision de produits finis ne
  fournit pas directement la quantité de matières à commander ;
  \item \textbf{faible traçabilité} : il est difficile de savoir quel modèle,
  quelle recette ou quel ajustement a conduit à une valeur donnée.
\end{itemize}

\subsection{Contraintes identifiées}

Les contraintes métier ont orienté l'architecture. Les données commerciales
doivent rester sur la machine de l'entreprise. Le système doit fonctionner sur un
poste Windows, supporter une mise à jour par fichiers lorsque la connexion SQL
n'est pas disponible, ne pas écraser l'historique si une extraction est vide ou
aberrante, et fournir des exports compréhensibles par des utilisateurs non
techniques. Enfin, les recettes restent en cours de fiabilisation ; le système doit
donc distinguer une recette exacte d'une estimation et rendre cette incertitude
visible.

\section{Solution proposée}

\subsection{Vue d'ensemble}

La solution retenue est une application locale en Python composée d'une chaîne de
traitement et d'un tableau de bord web. Elle conserve les données métier dans des
fichiers CSV et JSON lisibles, tout en proposant un accès automatisé en lecture à
la base SQL de la caisse. Elle produit deux niveaux de prévision complémentaires :

\begin{itemize}
  \item un système \textbf{mensuel}, destiné au plan de production, au budget et
  au calcul MRP ;
  \item un système \textbf{journalier}, destiné à la préparation quotidienne et
  au suivi court terme.
\end{itemize}

La chaîne complète va de la vente réelle à la recommandation : préparation des
données, sélection du modèle, ajustements calendrier et commandes, estimation de
l'incertitude, explosion des nomenclatures, calcul des coûts, génération d'exports
et restitution dans le dashboard.

\subsection{Avantages de la solution retenue}

Cette solution présente plusieurs avantages. Elle est reproductible, car un même
jeu de données et une même configuration conduisent au même résultat. Elle est
traçable, car le modèle retenu par produit, la source des recettes et les
ajustements sont externalisés. Elle est évolutive, car les modules de prévision,
de nomenclature et d'interface peuvent être améliorés séparément. Elle est enfin
compatible avec le niveau d'équipement existant : un poste Windows et un
navigateur suffisent pour consulter l'application.

\subsection{Originalité du projet}

L'originalité ne réside pas dans l'emploi isolé d'un modèle statistique. Elle
provient de l'enchaînement de plusieurs mécanismes métier dans une même solution :

\begin{itemize}
  \item comparaison de dix candidats de prévision et choix différent selon le produit ;
  \item réconciliation entre la somme des produits et la prévision agrégée ;
  \item apprentissage de l'impact des matchs à partir des jours comparables ;
  \item neutralisation des pics assimilables à de grosses commandes dans
  l'apprentissage, puis ajout exact des commandes futures connues ;
  \item conversion récursive des produits semi-finis en matières de base ;
  \item assistant guidé local, sans modèle génératif ni appel réseau, dont chaque
  réponse provient des fichiers calculés ;
  \item articulation entre précision statistique, marge de sécurité et usage
  opérationnel.
\end{itemize}

\section{Démarche de développement}

Nous avons adopté une démarche itérative proche d'un cycle incrémental. Chaque
itération produit un résultat testable : un import fiable, une prévision, un
export, une nouvelle vue du dashboard ou une amélioration des recettes. Ce choix
est plus adapté qu'un cycle strictement séquentiel, car la qualité des données et
les retours métier modifient régulièrement les règles de calcul.

Le cycle d'une itération comprend cinq étapes : observer le besoin, isoler les
données concernées, implémenter une règle ou un modèle, vérifier par test et
backtest, puis présenter le résultat dans un format exploitable. Le dépôt Git
conserve les évolutions et les correctifs. La version synchronisée comprend 22
commits depuis l'import initial public du code, avec une attention particulière à
l'automatisation SQL et à la fiabilisation des recettes.

\section{Planning prévisionnel}

Le tableau~\ref{tab:planning} présente le découpage prévisionnel adopté. Les
semaines sont volontairement exprimées relativement au démarrage afin que le
planning reste cohérent avec les dates administratives du stage.

\begin{table}[H]
\centering
\caption{Planning prévisionnel sur douze semaines}
\label{tab:planning}
\scriptsize
\setlength{\tabcolsep}{2pt}
\begin{tabular}{p{4.4cm}*{12}{C{0.46cm}}}
\toprule
\textbf{Étape} & \textbf{S1}&\textbf{S2}&\textbf{S3}&\textbf{S4}&\textbf{S5}&\textbf{S6}&\textbf{S7}&\textbf{S8}&\textbf{S9}&\textbf{S10}&\textbf{S11}&\textbf{S12}\\
\midrule
Analyse du besoin et des sources
&\cellcolor{paulGold}&\cellcolor{paulGold}&&&&&&&&&&\\
Nettoyage et consolidation des ventes
&&\cellcolor{paulGold}&\cellcolor{paulGold}&&&&&&&&&\\
Prototype de prévision mensuelle
&&&\cellcolor{paulGold}&\cellcolor{paulGold}&&&&&&&&\\
Prévision journalière et événements
&&&&\cellcolor{paulGold}&\cellcolor{paulGold}&&&&&&&\\
Nomenclatures et besoins matières
&&&&&\cellcolor{paulGold}&\cellcolor{paulGold}&\cellcolor{paulGold}&&&&\\
Tableau de bord et exports
&&&&&&\cellcolor{paulGold}&\cellcolor{paulGold}&\cellcolor{paulGold}&&&\\
Automatisation et déploiement
&&&&&&&&\cellcolor{paulGold}&\cellcolor{paulGold}&&\\
Tests, validation et documentation
&&&&&&&&&\cellcolor{paulGold}&\cellcolor{paulGold}&\cellcolor{paulGold}\\
\bottomrule
\end{tabular}
\end{table}

\section{Conclusion}

Ce chapitre a montré que le besoin dépasse une simple visualisation de ventes.
Le problème consiste à relier des sources hétérogènes à une décision de production
et d'achat, en mesurant l'incertitude et en conservant la confidentialité. La
solution proposée répond à ce cadre par une architecture locale, modulaire et
itérative. Le chapitre suivant traduit ce cadre en besoins précis et en cas
d'utilisation.
